\section{Edition Tools}

Sometimes the user needs to modify the result of a segmentation algorithm, 
group several segmentations into one, or create a new segmentation. These kind
of modifications can be done with the edition tools, located in the respective
toolbar.\\

\begin{figure}[H]
\centering
\includegraphics{fig/EditorToolbar}
\caption{Edition toolbar.}
\end{figure}

The first is a multi-state button that combines all the manual edition tools,
particularly:\\

\vspace{0.1cm}
\begin{tabular}{| m{1.3cm} | m{12cm} |}
\hline
\textbf{Button} & \textbf{Description}\\
\hline
\includegraphics[width=0.7cm]{../../frontend/rsc/pencil2D} &
\textbf{2D Brush}: Use a circle to draw on current slice. Use Ctrl+Wheel to
change brush radius.\\
\hline
\includegraphics[width=0.7cm]{../../frontend/rsc/pencil3D} &
\textbf{3D Brush}: Use a sphere to draw on current and adjacent slices. Use
Ctrl+Wheel to change brush radius.\\
\hline
\includegraphics[width=0.7cm]{../../frontend/rsc/eraser2D} &
\textbf{Eraser}: Ctrl click while using any of the previous brushes to erase
instead of draw.\\
\includegraphics[width=0.7cm]{../../frontend/rsc/eraser3D} & \\
\hline
\includegraphics[width=0.7cm]{../../frontend/rsc/lasso} &
\textbf{Contour}: Use a free-form contour to draw on current slice.\\
\hline
\end{tabular}
\vspace{0.3cm}

To start creating a new segmentation select your tool and then with no other
segmentation being selected, start drawing on the 2D views. To modify an existing
segmentation just select your tool with the segmentation you need the modify 
selected in the segmentation explorer.\\
When holding the control key, the brush tool will erase instead of drawing.\\

\vspace{0.3cm}
\begin{bclogo}[couleur = yellow!33, logo= \bcbook]
{Tip} As a visual help, the brush border color will change from blue (creating a
new segmentation) to green (editing current segmentation). 
\end{bclogo}
\vspace{0.3cm}

The contour drawing, when started, will allow us to put nodes on the view using the left
mouse button. You can maintain the left button down to draw continuously over the slice.
While the contour is ``open'' nodes can be added anywhere of the slice, until the last point
of the contour is added using the right mouse button or, when adding a node with the left
button, the new segment intersects with the contour. In that case the contour will delete
unused nodes and will set the contour as ``closed''.\\
A closed contour can be modified: a node can be added when the mouse cursor is near the edge
of the contour, nodes can be translated if the user clicks and drags the node with the left
button, and can be deleted if the user holds down the shift key while left-clicking over a node.
Also a contour can be tranlated over the plane (click and drag inside the contour) or simply
erased completely (using the backspace key).\\
In every one of these states the mouse cursor will change to indicate the outcome of the action.
\vspace{0.3cm}

\begin{tabular}{| m{1.3cm} | m{12cm} |}
\hline
\textbf{Cursor} & \textbf{Description}\\
\hline
Hand & Translate whole contour\\
\hline
Pointing hand & Translate node\\
\hline
\includegraphics[width=0.7cm]{../../frontend/rsc/cross-plus} & Add a node in this position.\\
\hline
\includegraphics[width=0.7cm]{../../frontend/rsc/cross-minus} & Delete a node in this position (holding shift key).\\
\hline
\end{tabular}
\vspace{0.3cm}

Sometimes, it can be convenient to use different tools to segment different
parts of the same segmentation and then combine the results to create the final 
segmentation. 
\vspace{0.3cm}

\begin{tabular}{| m{1.3cm} | m{13.5cm} |}
\hline
\textbf{Button} & \textbf{Description}\\
\hline
\includegraphics[width=0.7cm]{../../frontend/rsc/add} &
\textbf{Combine}: Create a unique segmentation using all voxels from all
selected segmentations.\\
\hline
\includegraphics[width=0.7cm]{../../frontend/rsc/remove} &
\textbf{Substract}: Use the first selected input as final segmentation and erase
all voxels which belong to any of the rest of the segmentations.\\
\hline
\end{tabular}
\vspace{0.3cm}

The edition toolbar also provides some operations that affect the morphology of the 
selected semgentation. To apply them the only requirement is to have the segmentation
we want selected when clicking on the action.\\
\vspace{0.3cm}

\begin{tabular}{| m{1.3cm} | m{13.5cm} |}
\hline
\textbf{Button} & \textbf{Description}\\
\hline
\includegraphics[width=0.7cm]{../../frontend/rsc/erode} &
\textbf{Erode}: Gradually shrink the boundaries of selected segmentations by eroding the surface using a structuring element (default 3 voxels).\\
\hline
\includegraphics[width=0.7cm]{../../frontend/rsc/dilate} &
\textbf{Dilate}: Gradually enlarge the boundaries of selected segmentations by dilating the surface using a structuring element (default 3 voxels). \\
\hline
\includegraphics[width=0.7cm]{../../frontend/rsc/open} &
\textbf{Opening}: Structural removal of some voxels from the edges of the segmmentations. Contrary to erode, opening preserve those parts of the segmentation which have a similar shape to the structuring element, or that can completely contain the structuring element, a sphere of given radius (default 3 voxels).\\
\hline
\includegraphics[width=0.7cm]{../../frontend/rsc/close} &
\textbf{Closing}: Structural enlargement of voxels from the edges of the segmmentations. Contrary to dilate, closing preserve those cavities of the segmentation which have a similar shape to the structuring element, or that can completely contain the structuring element, a sphere of given radius (default 3 voxels).\\
\hline
\includegraphics[width=0.7cm]{../../frontend/rsc/fillHoles} &
\textbf{Fill Holes}: Fill all the voxels which are totally surrounded by
segmentation voxels.\\
\hline
\end{tabular}
